% ============================================================================
%  SECCIÓN 3.19: IMPLEMENTACIÓN DE FORMULARIOS
% ============================================================================

\section{Implementaci\'on de formularios}

La implementaci\'on de formularios es uno de los requisitos fundamentales de
este avance, pues permite al usuario interactuar con la aplicaci\'on mediante
la ingresa de datos. Para las pantallas de inicio de sesi\'on y registro se
utilizaron las librer\'ias \texttt{react-hook-form} y \texttt{@hookform/resolvers}
con validaci\'on \texttt{Zod}, que permiten gestionar el estado del formulario,
validar los campos en tiempo real y mostrar mensajes de error al usuario.

\subsection{Librer\'ias utilizadas}

\begin{table}[H]
  \centering
  \caption{Librer\'ias para gesti\'on de formularios}
  \contenidoConFuente{%
    \begin{tablaAPA}
      \begin{tabular}{@{}p{3.4cm}p{4.0cm}p{8.2cm}@{}}
        \toprule
        \textbf{Librer\'ia} & \textbf{Versi\'on} & \textbf{Prop\'osito} \\
        \midrule
        \textit{react-hook-form} & \textit{\^{}7.85.0} & Gestor de formularios
        ligero que reduce re-renders innecesarios y simplifica la
        validaci\'on de campos. \\
        \addlinespace
        \textit{@hookform/resolvers} & \textit{\^{}5.9.0} & Integraci\'on de
        \textit{react-hook-form} con esquemas de validaci\'on externos como
        Zod. \\
        \addlinespace
        \textit{Zod} & \textit{\^{}4.4.3} & Librer\'ia de validaci\'on y
        definici\'on de esquemas con tipado est\'atico en TypeScript. \\
        \bottomrule
      \end{tabular}
    \end{tablaAPA}%
  }{Elaboraci\'on propia en base al archivo \texttt{package.json} del
  proyecto.}
\end{table}

\subsection{Definici\'on del esquema de validaci\'on}

Cada formulario se define mediante un esquema Zod que establece las reglas de
validaci\'on de cada campo. El esquema se compila est\'aticamente con
TypeScript, lo que permite inferir el tipo de datos del formulario
autom\'aticamente.

A continuaci\'on se presenta el esquema de validaci\'on del formulario de
registro:

\begin{lstlisting}
const registerSchema = z
  .object({
    fullName: z
      .string()
      .min(1, 'El nombre es obligatorio')
      .min(3, 'El nombre debe tener al menos 3 caracteres')
      .max(100, 'El nombre es demasiado largo'),
    email: z
      .string()
      .min(1, 'El email es obligatorio')
      .email('Ingresa un email valido'),
    password: z
      .string()
      .min(1, 'La contraseña es obligatoria')
      .min(6, 'La contraseña debe tener al menos 6 caracteres'),
    confirmPassword: z
      .string()
      .min(1, 'Confirma tu contraseña'),
  })
  .refine((data) => data.password === data.confirmPassword, {
    message: 'Las contraseñas no coinciden',
    path: ['confirmPassword'],
  });
\end{lstlisting}

El esquema del formulario de inicio de sesi\'on es m\'as sencillo, requiriendo
\'unicamente el campo de email y contrase\~na:

\begin{lstlisting}
const loginSchema = z.object({
  email: z
    .string()
    .min(1, 'El email es obligatorio')
    .email('Ingresa un email valido'),
  password: z
    .string()
    .min(1, 'La contraseña es obligatoria')
    .min(6, 'La contraseña debe tener al menos 6 caracteres'),
});
\end{lstlisting}

\subsection{Integraci\'on con react-hook-form}

El componente \texttt{useForm} de \texttt{react-hook-form} se conecta con el
esquema Zod mediante el \texttt{zodResolver}, que traduce las reglas de
validaci\'on de Zod al formato de react-hook-form. El hook \texttt{Controller}
se utiliza para controlar cada campo de texto individualmente.

\begin{lstlisting}
const {
  control,
  handleSubmit,
  formState: { errors, isSubmitting },
} = useForm<RegisterFormData>({
  resolver: zodResolver(registerSchema),
  defaultValues: {
    fullName: '',
    email: '',
    password: '',
    confirmPassword: '',
  },
});
\end{lstlisting}

Cada campo de texto se renderiza con el componente \texttt{Controller}, que
sincroniza el valor del campo con el estado del formulario:

\begin{lstlisting}
<Controller
  control={control}
  name="email"
  render={({ field: { onChange, onBlur, value } }) => (
    <TextInput
      style={[styles.input, errors.email && styles.inputError]}
      placeholder="tu@email.com"
      placeholderTextColor={colors.textSecondary}
      value={value}
      onChangeText={onChange}
      onBlur={onBlur}
      keyboardType="email-address"
      autoCapitalize="none"
    />
  )}
/>
{errors.email && (
  <Text style={styles.errorText}>{errors.email.message}</Text>
)}
\end{lstlisting}

\subsection{Env\'io del formulario}

El env\'io del formulario se gestiona mediante la funci\'on \texttt{onSubmit},
que recibe los datos validados y ejecuta la acci\'on correspondiente (inicio de
sesi\'on o registro) a trav\'es del contexto de autenticaci\'on. Durante el
env\'io, se muestra un indicador de carga (\texttt{ActivityIndicator}) y se
deshabilita el bot\'on para evitar env\'ios duplicados.

\begin{lstlisting}
const onSubmit = async (data: RegisterFormData) => {
  setServerError(null);
  setSuccessMessage(null);

  const { error } = await signUp(
    data.email, data.password, data.fullName
  );

  if (error) {
    setServerError(error);
    return;
  }

  setSuccessMessage(
    'Cuenta creada. Revisa tu email para confirmar tu cuenta.'
  );

  setTimeout(() => {
    router.replace('/auth/login');
  }, 2500);
};
\end{lstlisting}

\subsection{Manejo de errores del servidor}

Adem\'as de las validaciones locales del cliente, la aplicaci\'on gestiona
errores provenientes del servidor (Supabase), como credenciales inv\'alidas o
correos ya registrados. Estos errores se muestran en una caja de estilo
especial con fondo rojo claro, diferenci\'andolos de las validaciones de campo.

\begin{lstlisting}
{serverError && (
  <View style={styles.serverErrorBox}>
    <Text style={styles.serverErrorText}>{serverError}</Text>
  </View>
)}
\end{lstlisting}
